\documentclass[11pt]{article}
\usepackage[margin=1in]{geometry}
\usepackage[T1]{fontenc}
\usepackage[utf8]{inputenc}
\usepackage{longtable}
\usepackage{array}
\usepackage{hyperref}
\usepackage{enumitem}
\title{RQ1: General Modeling of Nanosatellite Missions and EPS}
\date{2026-04-10}

\begin{document}
\maketitle

\section{Research question}
How can nanosatellite missions and electrical power systems (EPS) be modeled in a general manner with minimal knowledge about their behavior?

\section{Bibliometric overview}

\begin{table}[h]
\centering
\begin{tabular}{lrr}
\hline
Year & Extracted papers \\
\hline
2013 & 2 \\
2014 & 2 \\
2016 & 1 \\
2017 & 1 \\
2021 & 1 \\
2022 & 2 \\
2023 & 2 \\
2024 & 2 \\
2025 & 2 \\
\hline
\end{tabular}
\caption{Publication-year distribution.}
\end{table}



\begin{table}[h]
\centering
\begin{tabular}{lrr}
\hline
Venue type & Extracted papers \\
\hline
Conference & 12 \\
Journal & 3 \\
\hline
\end{tabular}
\caption{Venue-type distribution.}
\end{table}

\section{Direct-study synthesis}

\begin{longtable}{|p{0.6cm}|p{0.8cm}|p{3.1cm}|p{4.0cm}|p{3.3cm}|p{3.3cm}|}
\hline
ID & Year & Study & Modeling approach & Main genericity mechanism(s) & Evidence relevant to RQ1 \\ \hline
\endfirsthead
\hline
ID & Year & Study & Modeling approach & Main genericity mechanism(s) & Evidence relevant to RQ1 \\ \hline
\endhead
S1 & 2013 & Model Based Systems Engineering (MBSE) Applied to Radio Aurora Explorer (RAX) CubeSat Mission Operational Scenarios & Standards-based SysML MBSE with structural, behavioral, state-machine, activity, and parametric models, integrated with external simulation/analysis tools (MagicDraw/Cameo Simulation Toolkit, ParaMagic, STK, MATLAB, PHX ModelCenter). & CubeSat meta-model describing CubeSat-specific concepts; domain-specific SysML extension for CubeSat modeling; reusable modeling framework covering mission, flight system, ground system, environment, subsystems, interfaces, behaviors, constraints, and measures of effectiveness & The integrated MBSE environment can represent realistic CubeSat operational scenarios and support analysis of both physical and functional states.; A single parametric model can be reused for multiple analysis causalities. \\ \hline
S2 & 2013 & Model-based software engineering for an imaging CubeSat and its extrapolation to other missions & OPEN-SBM, a service-oriented model-based software engineering approach using SRD-style service blocks and finite-state-machine reasoning & service-oriented black-box services; subsystem decomposition into reusable service blocks; SRD graphical abstraction & service-oriented abstraction can organize CubeSat software into interacting subsystem services; the approach supports autonomous coordination of health, payload, communications, attitude control, and power management \\ \hline
S3 & 2014 & Enterprise modeling for CubeSats & Enterprise-oriented MBSE using SysML architecture framework, subsystem libraries, parametric cost-estimating relationships, integrated behavioral simulation, and PLM linkage & SysML library of common CubeSat elements; top-down requirements-to-subsystem flow-down; modular subsystem/component/interface/task decomposition & A reusable CubeSat framework can model mission, spacecraft, and ground elements together; Integrated SysML plus analytical tools enabled mission trade studies and requirements verification \\ \hline
S4 & 2014 & Rapid model-based inter-disciplinary design of a CubeSat mission & Inter-disciplinary model-based mission design architecture in Matlab using analytical subsystem/mission models and NSGA-II-style multi-objective genetic optimization. & Baseline design structure matrix across disciplines; Parameter-driven inputs and requirements; OTS component database for subsystem selection & Architecture rapidly identifies Pareto sets for CubeSat mission trades; Analytical visibility methods avoid costly numerical propagation \\ \hline
S5 & 2016 & SPARC-1: A joint US/Sweden multi-mission modular nanosatellite platform & Modular plug-and-play spacecraft architecture using SPA standards, xTEDS self-description, layered middleware, structural modularity, standard interfaces, and configurable subsystem composition around a 6U platform. & SPA standard interfaces; xTEDS self-describing components; Self-discovery/join middleware & A modular 6U CubeSat platform can combine miniaturization with plug-and-play subsystem integration; SPA/xTEDS middleware enables reusable black-box component integration \\ \hline
S6 & 2023 & Crafting CanSats: A Novel Modular Design Paradigm for Scientific CanSats & Modular platform design standard based on standardized mating rings, interstage connectors, shared power/data bus, distributed microcontrollers, and uniform module naming/data exchange conventions. & Standard mechanical mating ring and hole pattern; Unified 8-pin interstage power/data connector; Shared I2C communication bus & Multiple module types were successfully manufactured and integrated under one standard; Students produced working modules within a semester \\ \hline
S7 & 2023 & Development of a Nanosatellite System Modeling Architecture for EIRSAT-1 & SysML-based MBSE using CATIA Magic Systems of Systems Architect, guided primarily by INCOSE OOSEM, with model integration to external Python-based test scripts and Magic Model Analyst for verification-flow assessment. & Use of SysML as a standardized modeling language; Adoption of OOSEM as a reusable methodological framework; Single integrated model environment with traceability from requirements to verification & MBSE verification models can be integrated with actual test setups and external scripts for automated requirement verification.; The CATIA Magic requirements table updates automatically from test outcomes. \\ \hline
S8 & 2024 & FloripaSat-2: An Open-Source Platform for CubeSats & Modular open-source service-platform architecture for CubeSats, defining subsystem decomposition, interfaces, buses, and reusable hardware/firmware patterns for EPS, OBDH, and TTC. & Modular service-module decomposition; Standardized subsystem interfaces; Support for multiple payloads and missions & FloripaSat-2 is presented as a general-purpose open-source CubeSat platform; The same platform supports multiple missions with different payloads \\ \hline
S9 & 2024 & Open-Source Data Formalization through Model-Based Systems Engineering for Concurrent Preliminary Design of CubeSats & MBSE formalization using a generic SysML CubeSat model augmented with UML stereotypes, plus a parser that extracts model data into a database for concurrent preliminary-design analyses. & Generalized CubeSat SysML model; Reusable UML stereotypes for subsystems and mission data; Generalized operating-mode schema with common activation criteria & A generic CubeSat model can serve as a central data source across preliminary design; UML stereotypes enable structured characterization of mission and subsystem data \\ \hline
S10 & 2025 & Concurrent engineering process for electrical architecture of nanosatellites preliminary phases & Model-based concurrent engineering workflow integrating system models (IDM-CIC), mission and orbit scenario models (Simu-CIC), power-profile generation (IDM-Power), and EPS energy-balance simulation (Opalis with Flows Manager) via standardized file interfaces. & Common technical reference and shared database across subsystem experts; Standardized file formats for interoperability between tools; Separation of concerns between system description, mission scenario, power profile, and EPS simulation & The interoperable workflow enabled rapid concurrent EPS analyses and trade-offs for a 16U agile nanosatellite.; For the studied mission, a back solar-panel configuration was preferred over a side configuration. \\ \hline
\end{longtable}

\end{document}
